% from ega2-1-fr.tex
\oldpage[II]{98}

de c') que l'homomorphisme correspondant
$\Gamma(X,\mathcal{O}_X^n)\to\Gamma(X,\mathcal{O}_X)$ est surjectif, ce qui
démontre b).

Enfin, a) implique d) (\hyperref[I.5.1.9.2-fr]{I, 5.1.9.2}) et d) implique
trivialement d'). Reste à montrer que d') entraîne c'). Or, si
$\mathcal{F}'$ est un sous-$\mathcal{O}_X$-Module quasi-cohérent de
$\mathcal{O}_X^n$, la filtration
$0\subset\mathcal{O}_X\subset\mathcal{O}_X^2\subset\cdots\subset
\mathcal{O}_X^n$ définit sur $\mathcal{F}'$ une filtration formée des
$\mathcal{F}'_k=\mathcal{O}_X^k\cap\mathcal{F}'$ ($0\leq k\leq n$), qui
sont des $\mathcal{O}_X$-Modules quasi-cohérents
(\hyperref[I.4.1.1-fr]{I, 4.1.1}), et
$\mathcal{F}'_{k+1}/\mathcal{F}'_k$ est isomorphe à un
sous-$\mathcal{O}_X$-Module quasi-cohérent de
$\mathcal{O}_X^{k+1}/\mathcal{O}_X^k=\mathcal{O}_X$, c'est-à-dire à un
faisceau quasi-cohérent d'idéaux de $\mathcal{O}_X$. L'hypothèse d')
entraîne donc
$\mathrm{H}^1(X,\mathcal{F}'_{k+1}/\mathcal{F}'_k)=0$~; la suite exacte de
cohomologie
$\mathrm{H}^1(X,\mathcal{F}'_k)\to\mathrm{H}^1(X,\mathcal{F}'_{k+1})
\to\mathrm{H}^1(X,\mathcal{F}'_{k+1}/\mathcal{F}'_k)=0$ permet alors de
prouver par récurrence sur $k$ que
$\mathrm{H}^1(X,\mathcal{F}'_k)=0$ pour tout $k$. C.Q.F.D.

\begin{remark}[5.2.1.1]
\label{II.5.2.1.1-fr}
Lorsque $X$ est un préschéma noethérien, on peut dans l'énoncé de c') et de
d'), remplacer «~quasi-cohérent~» par «~cohérent~». En effet, dans la
démonstration du fait que c') entraîne b), $\mathcal{J}$ et $\mathcal{J}'$
sont alors des faisceaux cohérents d'idéaux, et, d'autre part, tout
sous-Module quasi-cohérent d'un Module cohérent est cohérent
(\hyperref[I.6.1.1-fr]{I, 6.1.1})~; d'où la conclusion.
\end{remark}

\begin{corollary}[5.2.2]
\label{II.5.2.2-fr}
Soit $f:X\to Y$ un morphisme séparé quasi-compact. Les conditions suivantes
sont équivalentes~:
\begin{enumerate}
  \item[\rm{a)}] $f$ est un morphisme affine.
  \item[\rm{b)}] Le foncteur $f_*$ est exact dans la catégorie des
  $\mathcal{O}_X$-Modules quasi-cohérents.
  \item[\rm{c)}] Pour tout $\mathcal{O}_X$-Module quasi-cohérent
  $\mathcal{F}$, on a $\mathrm{R}^1f_*(\mathcal{F})=0$.
  \item[\rm{c')}] Pour tout faisceau d'idéaux quasi-cohérent $\mathcal{J}$
  de $\mathcal{O}_X$, on a $\mathrm{R}^1f_*(\mathcal{J})=0$.
\end{enumerate}
\end{corollary}

Toutes ces conditions étant locales sur $Y$, par définition du foncteur
$\mathrm{R}^1f_*$ (T, 3.7.3), on peut supposer que $Y$ est affine. Si $f$
est affine, $X$ est alors affine et la propriété b) n'est autre que
(\hyperref[I.1.6.4-fr]{I, 1.6.4}). Inversement, montrons que b) implique
a)~: pour tout $\mathcal{O}_X$-Module quasi-cohérent $\mathcal{F}$,
$f_*(\mathcal{F})$ est un $\mathcal{O}_Y$-Module quasi-cohérent
(\hyperref[I.9.2.2-fr]{I, 9.2.2, a)}). Par hypothèse le foncteur
$f_*(\mathcal{F})$ est exact en $\mathcal{F}$, et le foncteur
$\Gamma(Y,\mathcal{G})$ est exact en $\mathcal{G}$ (dans la catégorie des
$\mathcal{O}_Y$-Modules quasi-cohérents) puisque $Y$ est affine
(\hyperref[I.1.3.11-fr]{I, 1.3.11})~; donc
$\Gamma(Y,f_*(\mathcal{F}))=\Gamma(X,\mathcal{F})$ est exact en
$\mathcal{F}$, ce qui prouve notre assertion en vertu de
(\hyperref[II.5.2.1-fr]{5.2.1, c)}).

Si $f$ est affine, $f^{-1}(U)$ est affine pour tout ouvert affine $U$ de
$Y$ (\hyperref[II.1.3.2-fr]{1.3.2}), donc
$\mathrm{H}^1(f^{-1}(U),\mathcal{F})=0$
(\hyperref[II.5.2.1-fr]{5.2.1, d)}), ce qui par définition entraîne
$\mathrm{R}^1f_*(\mathcal{F})=0$. Supposons enfin vérifiée la condition
c')~; la suite exacte des termes de bas degré dans la suite spectrale de
Leray (G, II, 4.17.1 et I, 4.5.1) donne en particulier la suite exacte
\[
  0\to\mathrm{H}^1(Y,f_*(\mathcal{J}))\to
  \mathrm{H}^1(X,\mathcal{J})\to
  \mathrm{H}^0(Y,\mathrm{R}^1f_*(\mathcal{J})).
\]
Comme $Y$ est affine et $f_*(\mathcal{J})$ quasi-cohérent
(\hyperref[I.9.2.2-fr]{I, 9.2.2, a)}), on a
$\mathrm{H}^1(Y,f_*(\mathcal{J}))=0$
(\hyperref[II.5.2.1-fr]{5.2.1})~; l'hypothèse c') entraîne donc
$\mathrm{H}^1(X,\mathcal{J})=0$, et on conclut par
(\hyperref[II.5.2.1-fr]{5.2.1}) que $X$ est un schéma affine.

\begin{corollary}[5.2.3]
\label{II.5.2.3-fr}
Si $f:X\to Y$ est un morphisme affine, alors, pour tout
$\mathcal{O}_X$-Module quasi-cohérent $\mathcal{F}$, l'homomorphisme
canonique
$\mathrm{H}^1(Y,f_*(\mathcal{F}))\to\mathrm{H}^1(X,\mathcal{F})$ est
bijectif.
\end{corollary}

